% !TEX TS-program = xelatex
% !TEX encoding = UTF-8 Unicode
\documentclass[a4paper,12pt]{article}

% Поддержка Unicode и шрифтов для кроссплатформенности
\usepackage{fontspec}
\usepackage{polyglossia}
\usepackage{amsmath}

% Основные шрифты (кроссплатформенные)
\setmainfont{DejaVu Serif}
\setsansfont{DejaVu Sans}
\setmonofont{DejaVu Sans Mono}

% Для polyglossia кириллица
\newfontfamily\cyrillicfont{DejaVu Serif}
\newfontfamily\cyrillicfontsf{DejaVu Sans}
\newfontfamily\cyrillicfonttt{DejaVu Sans Mono}

% Настройки русского языка
\setdefaultlanguage{russian}
\setotherlanguages{english}
\defaultfontfeatures{Ligatures=TeX}

% Геометрия страницы
\usepackage[a4paper, margin=2cm]{geometry}

% Цвета для подсветки кода
\usepackage{xcolor}
\definecolor{codegreen}{rgb}{0,0.6,0}
\definecolor{codegray}{rgb}{0.5,0.5,0.5}
\definecolor{codepurple}{rgb}{0.58,0,0.82}
\definecolor{backcolour}{rgb}{0.95,0.95,0.92}

% Подсветка синтаксиса для Java и YAML
\usepackage{listings}
\lstdefinestyle{java}{
    backgroundcolor=\color{backcolour},
    commentstyle=\color{codegreen},
    keywordstyle=\color{magenta},
    numberstyle=\tiny\color{codegray},
    stringstyle=\color{codepurple},
    basicstyle=\ttfamily\footnotesize,
    breakatwhitespace=false,         
    breaklines=true,                 
    captionpos=b,                    
    keepspaces=true,                 
    numbers=left,                    
    numbersep=5pt,                  
    showspaces=false,                
    showstringspaces=false,
    showtabs=false,                  
    tabsize=2,
    frame=single,
    framesep=3pt,
    rulecolor=\color{gray},
    xleftmargin=10pt,
    framexleftmargin=10pt,
    framexrightmargin=5pt,
    framextopmargin=2pt,
    framexbottommargin=2pt
}

\lstdefinestyle{yaml}{
    backgroundcolor=\color{backcolour},
    commentstyle=\color{codegreen},
    keywordstyle=\color{magenta},
    numberstyle=\tiny\color{codegray},
    stringstyle=\color{codepurple},
    basicstyle=\ttfamily\footnotesize,
    breakatwhitespace=false,         
    breaklines=true,                 
    captionpos=b,                    
    keepspaces=true,                 
    numbers=left,                    
    numbersep=5pt,                  
    showspaces=false,                
    showstringspaces=false,
    showtabs=false,                  
    tabsize=2,
    frame=single,
    framesep=3pt,
    rulecolor=\color{gray},
    xleftmargin=10pt
}

\lstset{style=java}

% Заголовки разделов
\usepackage{titlesec}
\titleformat{\section}{\normalfont\Large\bfseries}{\thesection}{1em}{}
\titleformat{\subsection}{\normalfont\large\bfseries}{\thesubsection}{1em}{}

% Интервалы
\usepackage{setspace}
\onehalfspacing

% Гиперссылки
\usepackage{hyperref}
\hypersetup{
    colorlinks=true,
    linkcolor=blue,
    filecolor=magenta,      
    urlcolor=cyan,
}

% Названия разделов на русском
\addto\captionsrussian{\renewcommand{\contentsname}{Содержание}}
\addto\captionsrussian{\renewcommand{\listfigurename}{Список рисунков}}
\addto\captionsrussian{\renewcommand{\listtablename}{Список таблиц}}
\addto\captionsrussian{\renewcommand{\lstlistingname}{Листинг}}

% Метаданные документа
\title{Интеграция Apache Kafka в Spring Boot:\\Практическое руководство с KRaft и Docker}
\author{Конспект по разработке распределённых систем}
\date{\today}

\begin{document}

\maketitle
\tableofcontents

\section{Введение}
Данная лекция посвящена интеграции Apache Kafka в существующее Spring Boot-приложение для отработки навыков работы с event-driven архитектурой. 

\textbf{Ключевые особенности подхода:}
\begin{enumerate}
    \item \textbf{Kafka 4.1+ в режиме KRaft} — без ZooKeeper, с нативным консенсусом на базе Raft
    \item \textbf{Spring Kafka 4.0.x} — официальная поддержка современных версий Kafka
    \item \textbf{Ручная реализация} продюсера и консьюмера с использованием Spring Kafka (без высокоуровневых абстракций)
    \item \textbf{Логирование в файл} на чистом Java I/O в консьюмере (для отладки и обучения)
    \item \textbf{Docker Compose} для развёртывания Kafka и Spring Boot в контейнерах
    \item \textbf{IntelliJ IDEA} как среда разработки с настройками для работы с контейнерами
\end{enumerate}

\section{Архитектура решения}
\begin{figure}[h]
\centering
\fbox{\parbox{0.9\textwidth}{
\centering
\textbf{Схема взаимодействия:}\\[1em]
Spring Boot App $\xrightarrow{\text{KafkaTemplate}}$ Kafka Broker (KRaft) $\xrightarrow{\text{@KafkaListener}}$ File Logger\\[1em]
\textit{Все компоненты работают в Docker-контейнерах}
}}
\caption{Архитектура интеграции Kafka}
\end{figure}

\section{Настройка Docker Compose для Kafka в режиме KRaft}
Начиная с версии 4.0, Apache Kafka полностью перешёл на режим KRaft, исключив зависимость от ZooKeeper. Для локальной разработки используем \texttt{bitnami/kafka:latest} или \texttt{apache/kafka:latest}.

\subsection{docker-compose.yml}
\begin{lstlisting}[style=yaml,caption=Docker Compose для Kafka (KRaft mode)]
version: '3.8'

services:
  kafka:
    image: 'bitnami/kafka:latest'  # или 'apache/kafka:4.1.0'
    container_name: kafka-broker
    ports:
      - '9092:9092'   # для внутренних подключений (из других контейнеров)
      - '9094:9094'   # для внешних подключений (с хоста, IntelliJ IDEA)
    environment:
      # KRaft: режим процесса (контроллер + брокер для dev)
      - KAFKA_CFG_NODE_ID=0
      - KAFKA_CFG_PROCESS_ROLES=controller,broker
      # Настройка listeners
      - KAFKA_CFG_LISTENERS=PLAINTEXT://:9092,CONTROLLER://:9093,EXTERNAL://:9094
      - KAFKA_CFG_ADVERTISED_LISTENERS=PLAINTEXT://kafka:9092,EXTERNAL://localhost:9094
      - KAFKA_CFG_LISTENER_SECURITY_PROTOCOL_MAP=CONTROLLER:PLAINTEXT,PLAINTEXT:PLAINTEXT,EXTERNAL:PLAINTEXT
      # Кворум контроллеров (для KRaft)
      - KAFKA_CFG_CONTROLLER_QUORUM_VOTERS=0@kafka:9093
      - KAFKA_CFG_CONTROLLER_LISTENER_NAMES=CONTROLLER
      # Для dev-среды: один реплика-фактор
      - KAFKA_CFG_OFFSETS_TOPIC_REPLICATION_FACTOR=1
      - KAFKA_CFG_TRANSACTION_STATE_LOG_REPLICATION_FACTOR=1
      - KAFKA_CFG_TRANSACTION_STATE_LOG_MIN_ISR=1
      # Разрешить plaintext (только для разработки!)
      - ALLOW_PLAINTEXT_LISTENER=yes
      # Логирование
      - KAFKA_CFG_LOG4J_LOGGERS=kafka.controller=INFO,kafka.producer.async.DefaultEventHandler=INFO
    volumes:
      - kafka_data:/bitnami/kafka
    healthcheck:
      test: ["CMD", "kafka-broker-api-versions.sh", "--bootstrap-server", "localhost:9092"]
      interval: 10s
      timeout: 10s
      retries: 5
      start_period: 30s

  # Инициализация топиков при старте
  init-kafka:
    image: 'bitnami/kafka:latest'
    depends_on:
      kafka:
        condition: service_healthy
    entrypoint: ['/bin/sh', '-c']
    command: |
      "
      echo 'Ожидание доступности Kafka...'
      until kafka-topics.sh --bootstrap-server kafka:9092 --list > /dev/null 2>&1; do
        sleep 2
      done
      echo 'Создание топиков...'
      kafka-topics.sh --bootstrap-server kafka:9092 \
        --create --if-not-exists \
        --topic note-events \
        --replication-factor 1 --partitions 3
      kafka-topics.sh --bootstrap-server kafka:9092 \
        --create --if-not-exists \
        --topic audit-log \
        --replication-factor 1 --partitions 1
      echo 'Топики созданы:'
      kafka-topics.sh --bootstrap-server kafka:9092 --list
      "

  # Spring Boot приложение (собирается отдельно, см. раздел 4)
  spring-app:
    build:
      context: .
      dockerfile: Dockerfile
    container_name: notes-api
    ports:
      - '8080:8080'
    environment:
      - SPRING_PROFILES_ACTIVE=docker
      - KAFKA_BOOTSTRAP_SERVERS=kafka:9092
    depends_on:
      kafka:
        condition: service_healthy
    restart: unless-stopped

volumes:
  kafka_data:
\end{lstlisting}

\textbf{Важные замечания по конфигурации:}
\begin{enumerate}
    \item \texttt{KAFKA\_CFG\_PROCESS\_ROLES=controller,broker} — комбинированный режим для разработки; в production рекомендуется разделять роли
    \item \texttt{EXTERNAL} listener на порту 9094 позволяет подключаться к брокеру с хоста (например, из IntelliJ IDEA для отладки)
    \item \texttt{healthcheck} использует \texttt{kafka-broker-api-versions.sh} для надёжной проверки готовности брокера
    \item Для продакшена замените \texttt{ALLOW\_PLAINTEXT\_LISTENER=yes} на SSL/SASL аутентификацию
\end{enumerate}

\section{Настройка Spring Boot приложения}
\subsection{Зависимости (build.gradle.kts)}
\begin{lstlisting}[language=Java,caption=Добавление Spring Kafka]
plugins {
    id("org.springframework.boot") version "3.3.5"
    id("io.spring.dependency-management") version "1.1.6"
    kotlin("jvm") version "1.9.25" // опционально
    kotlin("plugin.spring") version "1.9.25"
}

dependencies {
    // Существующие зависимости...
    implementation("org.springframework.boot:spring-boot-starter-web")
    implementation("org.springframework.boot:spring-boot-starter-data-jpa")
    
    // Kafka integration
    implementation("org.springframework.kafka:spring-kafka:4.0.5") // актуальная версия
    implementation("io.jsonwebtoken:jjwt-api:0.13.0")
    implementation("io.jsonwebtoken:jjwt-impl:0.13.0")
    implementation("io.jsonwebtoken:jjwt-jackson:0.13.0")
    
    // Lombok, validation и т.д.
    compileOnly("org.projectlombok:lombok")
    annotationProcessor("org.projectlombok:lombok")
}
\end{lstlisting}

\subsection{Конфигурация application-docker.properties}
\begin{lstlisting}[style=java,caption=Настройки Kafka для Docker-среды]
# Kafka configuration
spring.kafka.bootstrap-servers=${KAFKA_BOOTSTRAP_SERVERS:kafka:9092}
spring.kafka.consumer.group-id=notes-service-group
spring.kafka.consumer.auto-offset-reset=earliest
spring.kafka.consumer.key-deserializer=org.apache.kafka.common.serialization.StringDeserializer
spring.kafka.consumer.value-deserializer=org.springframework.kafka.support.serializer.JsonDeserializer
spring.kafka.consumer.properties.spring.json.trusted.packages=com.example.demo.model,com.example.demo.dto

spring.kafka.producer.key-serializer=org.apache.kafka.common.serialization.StringSerializer
spring.kafka.producer.value-serializer=org.springframework.kafka.support.serializer.JsonSerializer

# Topic names
kafka.topic.note-events=note-events
kafka.topic.audit-log=audit-log

# Logging for Kafka (опционально, для отладки)
logging.level.org.apache.kafka=INFO
logging.level.org.springframework.kafka=WARN
\end{lstlisting}

\section{Ручная реализация Kafka Producer}
Используем \texttt{KafkaTemplate} из Spring Kafka, но реализуем логику отправки вручную (без @EventListener и т.п.).

\subsection{NoteEvent DTO}
\begin{lstlisting}[language=Java,caption=Событие для отправки в Kafka]
package com.example.demo.dto.kafka;

import lombok.*;
import java.time.Instant;

@Data
@Builder
@NoArgsConstructor
@AllArgsConstructor
public class NoteEvent {
    private String eventId;          // UUID для идемпотентности
    private Long noteId;
    private String action;           // "CREATED", "UPDATED", "DELETED"
    private String title;
    private Instant timestamp;
    private Long userId;         
}
\end{lstlisting}

\subsection{KafkaProducerService}
\begin{lstlisting}[language=Java,caption=Сервис отправки событий]
package com.example.demo.service.kafka;

import com.example.demo.dto.kafka.NoteEvent;
import lombok.RequiredArgsConstructor;
import lombok.extern.slf4j.Slf4j;
import org.apache.kafka.clients.producer.RecordMetadata;
import org.springframework.kafka.core.KafkaTemplate;
import org.springframework.kafka.support.SendResult;
import org.springframework.stereotype.Service;

import java.util.concurrent.CompletableFuture;
import java.util.concurrent.ExecutionException;

@Service
@RequiredArgsConstructor
@Slf4j
public class KafkaProducerService {

    private final KafkaTemplate<String, NoteEvent> kafkaTemplate;
    private final String topicName; // inject from properties

    /**
     * Отправка события с обработкой результата (синхронно для простоты)
     * В production рекомендуется использовать асинхронный подход с колбэками
     */
    public boolean sendNoteEvent(NoteEvent event) {
        try {
            // Отправка в топик с ключом = noteId для гарантированного порядка по заметке
            CompletableFuture<SendResult<String, NoteEvent>> future = 
                kafkaTemplate.send(topicName, event.getNoteId().toString(), event);
            
            // Блокирующее ожидание (для учебного примера)
            SendResult<String, NoteEvent> result = future.get();
            RecordMetadata metadata = result.getRecordMetadata();
            
            log.info("Event sent: partition={}, offset={}, timestamp={}", 
                    metadata.partition(), metadata.offset(), metadata.timestamp());
            return true;
            
        } catch (InterruptedException | ExecutionException e) {
            log.error("Failed to send event: {}", event.getEventId(), e);
            Thread.currentThread().interrupt();
            return false;
        }
    }
    
    /**
     * Асинхронная отправка с колбэком (рекомендуемый подход)
     */
    public void sendNoteEventAsync(NoteEvent event, Runnable onSuccess, Runnable onError) {
        kafkaTemplate.send(topicName, event.getNoteId().toString(), event)
            .whenComplete((result, ex) -> {
                if (ex == null) {
                    log.debug("Async send succeeded: offset={}", 
                        result.getRecordMetadata().offset());
                    if (onSuccess != null) onSuccess.run();
                } else {
                    log.error("Async send failed", ex);
                    if (onError != null) onError.run();
                }
            });
    }
}
\end{lstlisting}

\subsection{Интеграция с NoteController}
\begin{lstlisting}[language=Java,caption=Отправка событий из контроллера]
// В NoteController добавьте:
private final KafkaProducerService kafkaProducerService;

@PostMapping
public ResponseEntity<Note> createNote(@RequestBody @Valid NoteCreateDto noteDto) {
    // ... существующая логика сохранения ...
    
    // Отправка события после успешного сохранения
    var event = NoteEvent.builder()
        .eventId(UUID.randomUUID().toString())
        .noteId(saved.getId())
        .action("CREATED")
        .title(saved.getTitle())
        .timestamp(Instant.now())
        .userId(getCurrentUserId()) // из SecurityContextHolder
        .build();
    
    kafkaProducerService.sendNoteEventAsync(event, 
        () -> log.info("Event sent for note {}", saved.getId()),
        () -> log.warn("Failed to send event for note {}", saved.getId())
    );
    
    return new ResponseEntity<>(saved, HttpStatus.CREATED);
}

// Аналогично для updateNote и deleteNote с action="UPDATED"/"DELETED"
\end{lstlisting}

\section{Ручная реализация Kafka Consumer с логированием в файл}
Для отработки навыков реализуем консьюмер, который читает события и записывает их в файл \textbf{на чистом Java I/O} (без Logback/SLF4J для файла).

\subsection{KafkaConsumerConfig}
\begin{lstlisting}[language=Java,caption=Конфигурация консьюмера]
package com.example.demo.config.kafka;

import com.example.demo.dto.kafka.NoteEvent;
import org.apache.kafka.clients.consumer.ConsumerConfig;
import org.apache.kafka.common.serialization.StringDeserializer;
import org.springframework.beans.factory.annotation.Value;
import org.springframework.context.annotation.Bean;
import org.springframework.context.annotation.Configuration;
import org.springframework.kafka.config.ConcurrentKafkaListenerContainerFactory;
import org.springframework.kafka.core.ConsumerFactory;
import org.springframework.kafka.core.DefaultKafkaConsumerFactory;
import org.springframework.kafka.support.serializer.JsonDeserializer;

import java.util.HashMap;
import java.util.Map;

@Configuration
public class KafkaConsumerConfig {

    @Value("${spring.kafka.bootstrap-servers}")
    private String bootstrapServers;

    @Value("${spring.kafka.consumer.group-id}")
    private String groupId;

    @Bean
    public ConsumerFactory<String, NoteEvent> consumerFactory() {
        Map<String, Object> props = new HashMap<>();
        props.put(ConsumerConfig.BOOTSTRAP_SERVERS_CONFIG, bootstrapServers);
        props.put(ConsumerConfig.GROUP_ID_CONFIG, groupId);
        props.put(ConsumerConfig.AUTO_OFFSET_RESET_CONFIG, "earliest");
        props.put(ConsumerConfig.KEY_DESERIALIZER_CLASS_CONFIG, StringDeserializer.class);
        props.put(ConsumerConfig.VALUE_DESERIALIZER_CLASS_CONFIG, JsonDeserializer.class);
        props.put(JsonDeserializer.TRUSTED_PACKAGES, 
            "com.example.demo.dto.kafka,com.example.demo.model");
        props.put(JsonDeserializer.VALUE_DEFAULT_TYPE, NoteEvent.class);
        // Для надёжности в production:
        props.put(ConsumerConfig.ENABLE_AUTO_COMMIT_CONFIG, false);
        
        return new DefaultKafkaConsumerFactory<>(props);
    }

    @Bean
    public ConcurrentKafkaListenerContainerFactory<String, NoteEvent> 
            kafkaListenerContainerFactory() {
        ConcurrentKafkaListenerContainerFactory<String, NoteEvent> factory = 
            new ConcurrentKafkaListenerContainerFactory<>();
        factory.setConsumerFactory(consumerFactory());
        factory.setConcurrency(3); // параллельная обработка
        return factory;
    }
}
\end{lstlisting}

\subsection{NoteEventConsumer с ручным логированием в файл}
\begin{lstlisting}[language=Java,caption=Консьюмер с записью в файл через Java I/O]
package com.example.demo.consumer;

import com.example.demo.dto.kafka.NoteEvent;
import lombok.extern.slf4j.Slf4j;
import org.springframework.kafka.annotation.KafkaListener;
import org.springframework.kafka.support.Acknowledgment;
import org.springframework.stereotype.Component;

import java.io.BufferedWriter;
import java.io.IOException;
import java.nio.file.Files;
import java.nio.file.Path;
import java.nio.file.Paths;
import java.nio.file.StandardOpenOption;
import java.time.LocalDateTime;
import java.time.format.DateTimeFormatter;

@Component
@Slf4j
public class NoteEventConsumer {

    private static final Path LOG_FILE = Paths.get("kafka-events.log");
    private static final DateTimeFormatter DT_FORMATTER = 
        DateTimeFormatter.ofPattern("yyyy-MM-dd HH:mm:ss.SSS");

    /**
     * Обработчик событий из топика note-events
     * manual ack для контроля коммита после успешной записи в файл
     */
    @KafkaListener(topics = "${kafka.topic.note-events}", groupId = "${spring.kafka.consumer.group-id}")
    public void consumeNoteEvent(NoteEvent event, Acknowledgment ack) {
        try {
            // 1. Обработка события (бизнес-логика)
            processEvent(event);
            
            // 2. Ручная запись в файл на чистом Java I/O (требование задания)
            writeEventToFile(event);
            
            // 3. Подтверждение обработки (ручной коммит)
            ack.acknowledge();
            
        } catch (Exception e) {
            log.error("Error processing event: {}", event.getEventId(), e);
            // Не делаем ack — событие будет повторно доставлено
            // В production: добавить dead-letter queue
        }
    }

    private void processEvent(NoteEvent event) {
        // Пример бизнес-логики: обновление кэша, отправка уведомлений и т.п.
        log.info("Processing {} event for note #{}", event.getAction(), event.getNoteId());
    }

    /**
     * Запись события в файл с использованием стандартного Java I/O
     * (без использования логгеров для файла — как требовалось в задании)
     */
    private void writeEventToFile(NoteEvent event) throws IOException {
        String logEntry = String.format("[%s] EVENT: id=%s, noteId=%d, action=%s, title='%s', user=%s%n",
            LocalDateTime.now().format(DT_FORMATTER),
            event.getEventId(),
            event.getNoteId(),
            event.getAction(),
            event.getTitle(),
            event.getUserId()
        );
        
        // append = true, создаём файл если не существует
        try (BufferedWriter writer = Files.newBufferedWriter(
                LOG_FILE, 
                StandardOpenOption.CREATE, 
                StandardOpenOption.APPEND)) {
            writer.write(logEntry);
            // flush происходит автоматически при закрытии try-with-resources
        }
    }
}
\end{lstlisting}

\textbf{Пояснения:}
\begin{enumerate}
    \item \texttt{Acknowledgment ack} позволяет вручную подтвердить обработку сообщения (отключите \texttt{enable-auto-commit})
    \item \texttt{Files.newBufferedWriter} с \texttt{StandardOpenOption.APPEND} обеспечивает потокобезопасную запись
    \item В production добавьте ротацию логов (например, через \texttt{java.nio.file.attribute.FileAttribute} или внешний логгер)
    \item Для идемпотентности используйте \texttt{eventId} для фильтрации дубликатов
\end{enumerate}

\section{Dockerfile для Spring Boot приложения}
Используем multi-stage build для минимизации размера образа.

\begin{lstlisting}[style=yaml,caption=Dockerfile для Spring Boot]
# Stage 1: Build
FROM eclipse-temurin:21-jdk-jammy AS builder
WORKDIR /app

# Копируем только зависимости для кэширования слоя
COPY build.gradle.kts settings.gradle.kts ./
COPY gradle gradle
COPY gradlew ./
RUN chmod +x gradlew
RUN ./gradlew dependencies --no-daemon || true

# Копируем исходный код и собираем
COPY src ./src
RUN ./gradlew bootJar -x test --no-daemon

# Stage 2: Runtime
FROM eclipse-temurin:21-jre-jammy
WORKDIR /app

# Создаём non-root пользователя для безопасности
RUN groupadd -r spring && useradd -r -g spring spring
USER spring:spring

# Копируем JAR из builder stage
COPY --from=builder --chown=spring:spring /app/build/libs/*.jar app.jar

# Health check
HEALTHCHECK --interval=30s --timeout=3s --start-period=40s --retries=3 \
  CMD curl -f http://localhost:8080/actuator/health || exit 1

EXPOSE 8080
ENTRYPOINT ["java", "-XX:+UseContainerSupport", "-XX:MaxRAMPercentage=75.0", "-jar", "app.jar"]
\end{lstlisting}

\section{Настройка IntelliJ IDEA}
\subsection{Запуск в контейнерах}
\begin{enumerate}
    \item Установите плагин \textbf{Docker} (File → Settings → Plugins)
    \item Добавьте \texttt{docker-compose.yml} в Docker tool window
    \item Для отладки: запустите только Kafka-сервисы, а Spring Boot — локально с профилем \texttt{docker}:
    \begin{itemize}
        \item Run → Edit Configurations → Spring Boot → VM options: \texttt{-Dspring.profiles.active=docker}
        \item Environment variables: \texttt{KAFKA\_BOOTSTRAP\_SERVERS=localhost:9094} (EXTERNAL listener)
    \end{itemize}
    \item Для полного запуска в Docker: используйте конфигурацию \texttt{docker-compose up}
\end{enumerate}

\subsection{Отладка Kafka-потребителей}
\begin{itemize}
    \item Установите точку останова в \texttt{NoteEventConsumer.consumeNoteEvent()}
    \item Отправьте событие через REST API или Postman
    \item IDEA остановится в консьюмере — можно инспектировать \texttt{NoteEvent}
    \item Проверьте файл \texttt{kafka-events.log} в рабочей директории приложения
\end{itemize}

\section{Тестирование интеграции}
\subsection{Проверка через curl}
\begin{lstlisting}[style=yaml,caption=Тестовые запросы]
# 1. Регистрация и получение токена
curl -X POST http://localhost:8080/api/auth/register \
  -H "Content-Type: application/json" \
  -d '{"username":"test","password":"test1234"}'

curl -X POST http://localhost:8080/api/auth/login \
  -H "Content-Type: application/json" \
  -d '{"username":"test","password":"test1234"}'
# Ответ: {"token":"eyJhbGciOiJIUzI1NiJ9..."}

# 2. Создание заметки (триггерит Kafka-событие)
curl -X POST http://localhost:8080/api/notes \
  -H "Authorization: Bearer <token>" \
  -H "Content-Type: application/json" \
  -d '{"title":"Test Note","content":"Hello Kafka!"}'

# 3. Проверка файла с событиями
docker exec notes-api cat /app/kafka-events.log
# Или локально, если приложение запущено вне Docker:
tail -f kafka-events.log
\end{lstlisting}

\subsection{Мониторинг через Kafka CLI}
\begin{lstlisting}[style=yaml,caption=Полезные команды для отладки]
# Подключиться к контейнеру с инструментами:
docker-compose exec kafka bash

# Просмотр топиков:
kafka-topics.sh --bootstrap-server kafka:9092 --list

# Чтение сообщений из топика (с начала):
kafka-console-consumer.sh --bootstrap-server kafka:9092 \
  --topic note-events \
  --from-beginning \
  --property print.key=true \
  --property print.value=true

# Просмотр лага консьюмера:
kafka-consumer-groups.sh --bootstrap-server kafka:9092 \
  --group notes-service-group --describe
\end{lstlisting}

\section{Рекомендации для production}
\begin{enumerate}
    \item \textbf{Безопасность}: замените \texttt{PLAINTEXT} на \texttt{SASL\_SSL} с JWT/SASL аутентификацией
    \item \textbf{Масштабирование}: разделите роли контроллера и брокера, используйте 3+ ноды для кворума
    \item \textbf{Надёжность}: включите идемпотентный продюсер (\texttt{enable.idempotence=true}) и транзакции при необходимости
    \item \textbf{Мониторинг}: подключите Prometheus/Grafana через JMX-экспортер Kafka
    \item \textbf{Логирование}: замените ручную запись в файл на структурированные логи с ротацией (Logback + FileAppender)
    \item \textbf{Dead Letter Queue}: настройте \texttt{ErrorHandlingDeserializer} и топик для ошибок
\end{enumerate}

\section{Заключение}
Интеграция Kafka в Spring Boot с использованием современного стека (KRaft, Spring Kafka 4.x, Docker) позволяет:
\begin{itemize}
    \item Отработать навыки работы с event-driven архитектурой без избыточных абстракций
    \item Понять внутреннее устройство продюсеров/консьюмеров на практике
    \item Обеспечить переносимость среды разработки через Docker
    \item Подготовить основу для масштабирования и добавления сложных сценариев (CQRS, Saga, event sourcing)
\end{itemize}

Данный подход соответствует современным практикам: отказ от ZooKeeper в пользу KRaft, использование официальных образов Apache/Bitnami, и применение multi-stage Docker-сборки для минимизации attack surface.

\section*{Приложение: Полезные ссылки}
\begin{itemize}
    \item Официальная документация Kafka 4.1 (KRaft): \url{https://kafka.apache.org/41/operations/kraft/}
    \item Spring Kafka Reference: \url{https://docs.spring.io/spring-kafka/docs/4.0.x/reference/html/}
    \item Docker Compose для Kafka без ZooKeeper: \url{https://theexceptioncatcher.com/2025/06/kafka-in-2025-a-clean-docker-compose-setup-without-zookeeper/}
    \item Multi-stage builds для Spring Boot: \url{https://bjoernkw.com/2025/02/02/a-guide-to-docker-multi-stage-builds-for-spring-boot-by-catherine-edelveis/}
\end{itemize}

\end{document}
